\chapter*{Lời mở đầu}
\addcontentsline{toc}{chapter}{Lời mở đầu}

Một tiết học vui không phải là một tiết học toàn trò đùa. Nó là tiết học mà lời qua tiếng lại đủ thông minh để kéo học sinh nhập cuộc, đủ nhẹ để không căng, và đủ gọn để không phá nhịp dạy.

Cuốn sách này được viết cho đúng tinh thần đó. Mỗi mục là một màn đối đáp ngắn giữa thầy cô và học sinh, kèm theo cách dùng, biến thể nhanh và lưu ý về không khí. Mục tiêu không phải để thuộc lòng câu chữ, mà để người dạy có thêm cảm giác nhịp lớp: lúc nào nên đáp nhẹ, lúc nào nên dí dỏm, lúc nào nên dừng đúng chỗ.

Trong cuốn này, đối đáp được chia theo bốn tình huống gần nhất với lớp học thật:

\begin{itemize}
\item Mở tiết và kéo sự chú ý.
\item Gặp học sinh lì, lười, nói đỡ.
\item Góp ý mà không làm không khí nặng.
\item Dùng trong chủ nhiệm, sinh hoạt lớp, giờ cuối tuần.
\end{itemize}

Một câu đáp hay không làm lớp sợ hơn. Nó làm lớp tỉnh hơn, cười hơn, và chịu học hơn. Đó là điều cuốn sách này hướng tới.
